\section{Closed Gamma primitive certificates and analytic infinite jets}
\label{hp:section}

The newer of the two incoming \emph{Universal Gamma Normal Forms}
packages gives every canonical coefficient of the inverse-Gamma
differential hierarchy. Its research questions~1 and~6 ask, respectively,
for a nonrecursive formula for the compensating total derivative and
for analytic meaning of the infinite-jet generator. We answer those
questions here. The algebraic device used for the first answer is the
classical one-dimensional Euler homotopy operator
\cite{PooleHereman,HeremanDeconinckPoole}. Its application to the incoming
coefficient theorem, rather than the homotopy operator itself, is the
continuation result. The convergence theorem below supplies an analytic
domain for actual holomorphic logarithmic derivatives.

\subsection{Definitions and the inherited coefficient formula}

Let
\[
 \mathcal R=\mathbb Q[u_0,u_1,\ldots],\qquad Du_j=u_{j+1},
 \qquad \operatorname{wt}(u_j)=j+1.
\]
The factor degree of $u^{\boldsymbol a}=\prod_j u_j^{a_j}$ is
$|\boldsymbol a|=\sum_j a_j$. A monomial is canonical if it is a
power of $u_0$, or if its largest positive index occurs at least twice.
The canonical complement of $D\mathcal R$ gives unique polynomials
$\mathcal N(P),\mathcal H(P)$ satisfying
\begin{equation}\label{hp:decomposition}
 P=\mathcal N(P)+D\mathcal H(P),\qquad
 [1]\mathcal H(P)=0.
\end{equation}
This complement is classical \cite[Proposition~4.10]{Guo2014}.
For completeness, a unique largest derivative $u_j$ can be removed by
writing a suitable multiple of the monomial with $u_j$ replaced by
$u_{j-1}$ as a total derivative. The remainder has lower highest
index. This proves existence by termination. If $H$ is nonconstant
with highest index $j$, then $DH$ contains the nonzero part
$u_{j+1}\partial H/\partial u_j$, whose monomials are noncanonical.
Thus a canonical total derivative is zero, proving uniqueness.

Define
\begin{equation}\label{hp:hierarchy}
 A_0=1,\qquad A_{m+1}=u_0A_m-DA_m,
 \qquad N_m=\mathcal N(A_m),\quad H_m=\mathcal H(A_m).
\end{equation}
Their raw generating series is
\begin{equation}\label{hp:raw-generator}
 \sum_{m\ge0}A_m\frac{t^m}{m!}
 =\exp\left(\sum_{j\ge0}
                   \frac{(-1)^j u_jt^{j+1}}{(j+1)!}\right).
\end{equation}
Consequently, at weight $m=\sum(j+1)a_j$,
\begin{equation}\label{hp:raw-coefficient}
 a(\boldsymbol a):=[u^{\boldsymbol a}]A_m
 =\frac{(-1)^{m-|\boldsymbol a|}m!}
        {\prod_j\bigl(a_j!\,((j+1)!)^{a_j}\bigr)}.
\end{equation}

Here is the precise part of the incoming theorem used below
\cite{UniversalIncoming}. Set
\begin{equation}\label{hp:C}
 C_j(v)=\frac{(v+1)^{j+1}-(v-1)^{j+1}}
                  {2^{j+1}(j+1)!}.
\end{equation}
The coefficient $n(\boldsymbol a)=[u^{\boldsymbol a}]N_m$ is one
for $u_0^m$, zero for a unique highest positive derivative, and zero
when $\sum j a_j$ is odd. In the remaining case, with highest index
$d\ge1$ and $a_d\ge2$, it is
\begin{equation}\label{hp:normal-coefficient}
 n(\boldsymbol a)=
 \frac{(-1)^d m!}{4^d d!(d+1)!\prod_j a_j!}
 \int_0^1(1-v^2)^d
     C_d(v)^{a_d-2}\prod_{j=1}^{d-1}C_j(v)^{a_j}\,dv.
\end{equation}
The integrand is a polynomial, so this is an explicit finite rational
number. We retain attribution for this coefficient formula; the new
primitive formula starts with it as an input.

\subsection{A nonrecursive formula for every compensator}

For a differential polynomial define
\begin{align}
 E(P)&=\sum_{j\ge0}(-D)^j\frac{\partial P}{\partial u_j},
                                      \label{hp:Euler}\\
 B(P)&=\sum_{j\ge1}\sum_{k=0}^{j-1}
           u_{j-1-k}(-D)^k\frac{\partial P}{\partial u_j}.
                                      \label{hp:boundary}
\end{align}
Let $K$ denote the factor-degree operator:
$K(u^{\boldsymbol a})=|\boldsymbol a|u^{\boldsymbol a}$.
Its inverse is defined on polynomials without constant term by
division of each coefficient by its positive factor degree.

\begin{lemma}[Euler homotopy identity]\label{hp:homotopy-lemma}
For every $P\in\mathcal R$,
\begin{equation}\label{hp:Euler-boundary}
                 KP=u_0E(P)+DB(P).
\end{equation}
If $P$ has zero constant term and $E(P)=0$, then
\begin{equation}\label{hp:closed-homotopy}
                   D\bigl(BK^{-1}P\bigr)=P.
\end{equation}
\end{lemma}

\begin{proof}
For an arbitrary polynomial $Q$ and $j\ge1$, direct differentiation
of a telescoping sum gives
\[
 u_jQ-u_0(-D)^jQ
 =D\left(\sum_{k=0}^{j-1}u_{j-1-k}(-D)^kQ\right).
\]
Substitute $Q=\partial P/\partial u_j$ and sum over $j$, including
the $j=0$ term separately. Euler's identity for factor degree yields
\eqref{hp:Euler-boundary}. The operators $D$ and $B$ preserve factor
degree, while $E$ lowers it by one. Apply the identity separately
to each positive factor-degree component of $P$, then divide by that
degree. This proves \eqref{hp:closed-homotopy}. In particular no
analytic integration and no recursive integration by parts is involved.
\end{proof}

\begin{theorem}[Complete finite primitive formula]\label{hp:closed}
For every $m\ge1$,
\begin{equation}\label{hp:Hclosed}
 \boxed{\displaystyle H_m=BK^{-1}(A_m-N_m).}
\end{equation}
Equivalently, let $p(\boldsymbol a)=a(\boldsymbol a)-n(\boldsymbol a)$
with the explicit coefficients in
\eqref{hp:raw-coefficient}--\eqref{hp:normal-coefficient}. Then
\begin{equation}\label{hp:finite-sum}
 \boxed{\displaystyle
 H_m=\sum_{\sum(j+1)a_j=m}\frac{p(\boldsymbol a)}{|\boldsymbol a|}
       \sum_{j\ge1}a_j\sum_{k=0}^{j-1}(-1)^k
                  u_{j-1-k}D^k u^{\boldsymbol a-\boldsymbol e_j}.}
\end{equation}
All sums are finite. Together with \eqref{hp:normal-coefficient},
this determines both polynomials in the canonical decomposition
without any recursive normal reduction.
\end{theorem}

\begin{proof}
The polynomial $A_m-N_m$ is a total derivative by
\eqref{hp:decomposition}. The identity $E(DQ)=0$ follows by cancellation
from
$\partial_{u_j}D=D\partial_{u_j}+\partial_{u_{j-1}}$.
Lemma~\ref{hp:homotopy-lemma} therefore shows that the right side of
\eqref{hp:Hclosed} differentiates to $A_m-N_m$.
It has zero constant term, so uniqueness in
\eqref{hp:decomposition} identifies it with $H_m$.
Expanding the two operators gives \eqref{hp:finite-sum}.
\end{proof}

If a formula without powers of $D$ is desired, every term is expanded by
\begin{equation}\label{hp:Dexpansion}
 D^k u^{\boldsymbol b}
 =k!\sum_{\substack{\sum_{r\ge0}b_{i,r}=b_i\ (i\ge0)\\
                         \sum_{i,r}r b_{i,r}=k}}
       \prod_i\frac{b_i!}{\prod_r b_{i,r}!(r!)^{b_{i,r}}}
       \prod_{i,r}u_{i+r}^{b_{i,r}}.
\end{equation}
This is simply the multinomial Leibniz rule. Thus the word
``nonrecursive'' describes a literal finite coefficient formula,
not an instruction to run an integration-by-parts recurrence in disguise.

\subsection{An explicit higher Gamma primitive}

For example,
\begin{align}
 N_6={}&u_0^6-15u_0^2u_1^2+u_2^2,\label{hp:N6}\\
 H_6={}&-3u_0^5+20u_0^3u_1-15u_0^2u_2-15u_0u_1^2
             +6u_0u_3+9u_1u_2-u_4.\label{hp:H6}
\end{align}
At weight eight the normal part is
\begin{equation}\label{hp:N8}
 N_8=u_0^8-70u_0^4u_1^2+28u_0^2u_2^2-7u_1^4-u_3^2,
\end{equation}
and the companion polynomial, supplied completely by the same formula, is
\begin{align}
 H_8={}&-4u_0^7+56u_0^5u_1-70u_0^4u_2-140u_0^3u_1^2
              +56u_0^3u_3\notag\\
      &+252u_0^2u_1u_2-28u_0^2u_4+112u_0u_1^3
              -112u_0u_1u_3\notag\\
      &-84u_0u_2^2+8u_0u_5-98u_1^2u_2+20u_1u_4
              +36u_2u_3-u_6.\label{hp:H8}
\end{align}
These low-weight examples are illustrations of the closed construction,
not claims of new historical evaluations.

For a smooth function $U$, substitute $u_j=U^{(j)}$ and let $Y'=-UY$.
The inherited weighted identity now has a fully nonrecursive right side:
\begin{equation}\label{hp:weighted}
 \boxed{\displaystyle
 \int Y\bigl(N_m(U)+U H_m(U)\bigr)\,dx
       =-Y\bigl(A_{m-1}(U)+H_m(U)\bigr)+C.}
\end{equation}
Indeed $(YA_{m-1})'=-YA_m$ and
$(YH_m)'=Y(DH_m-UH_m)$. The formula follows on adding these two
identities. The compensator $UH_m$ and the endpoint value of $YH_m$
are part of the formula; neither can be discarded.

Taking $U=\psi$ and $Y=1/\Gamma$ gives exact Gamma--polygamma
antiderivatives at every weight. Taking instead
\begin{equation}\label{hp:StieltjesU}
 F_q(x)=\frac{(-1)^q}{q+1}\zeta^{(q+1)}(0,x),\quad
 U_q(x)=-\gamma_q(x),\quad Y_q(x)=e^{-F_q(x)}
\end{equation}
gives the corresponding all-index Stieltjes primitives.
Here $\zeta^{(q+1)}$ denotes an order derivative, not an argument
derivative. The standard identity
$\partial_x\zeta(s,x)=-s\zeta(s+1,x)$ and the Laurent convention
$\zeta(1+s,x)=s^{-1}+\sum_{q\ge0}(-1)^q\gamma_q(x)s^q/q!$
give $F_q'=-\gamma_q=U_q$ \cite{DLMFHurwitz}.
Thus $Y_q'=-U_qY_q$, as required. At $q=0$,
$F_0=\log\Gamma-\tfrac12\log(2\pi)$.

\subsection{The infinite-jet formula is analytic for holomorphic fields}

The incoming coefficient theorem is formal in infinitely many jets.
The following result permits their actual substitution from a
holomorphic function. Define the entire function
\[
 \phi(z)=\frac{e^z-1-z}{z^2},\qquad \phi(0)=\frac12.
\]
For $x$ in a domain of holomorphy of $U$, put
$x_j=t^{j+1}U^{(j)}(x)$ and
$S_{r-1}(v)=\sum_{j=1}^{r-1}C_j(v)x_j$.

\begin{theorem}[Analytic normal generator]\label{hp:analyticN}
Suppose $U$ is holomorphic on $|x-x_0|<R$. The series
\begin{equation}\label{hp:Nanalytic}
 \boxed{\displaystyle
 \mathfrak N_U(x,t)=e^{tU(x)}\left[
 1+\frac12\sum_{r\ge1}
 \frac{(-1)^r x_r^2}{4^r r!(r+1)!}
 \int_{-1}^1(1-v^2)^r e^{S_{r-1}(v)}
                         \phi(C_r(v)x_r)\,dv\right]}
\end{equation}
converges absolutely and locally uniformly when
\begin{equation}\label{hp:normaldomain}
                        |x-x_0|+|t|<R.
\end{equation}
It is holomorphic there and has Taylor series
\begin{equation}\label{hp:normalTaylor}
            \mathfrak N_U(x,t)=\sum_{m\ge0}N_m(U(x))\frac{t^m}{m!}.
\end{equation}
\end{theorem}

\begin{proof}
Cover any compact subset of \eqref{hp:normaldomain} by finitely
many sufficiently small compact polydisks. On each such polydisk
choose $\rho$ with $\sup|t|<\rho<R-\sup|x-x_0|$. Cauchy's
estimate supplies $M>0$ such that, uniformly there,
\[
 |U^{(r)}(x)|\le Mr!\rho^{-r},\qquad |t|<\rho.
\]
For real $-1\le v\le1$, the explicit positive-coefficient formula
for $C_j$ gives $|C_j(v)|\le1/(j+1)!$. Hence
\[
 \sum_{j\ge1}|C_j(v)x_j|
 \le M|t|\sum_{j\ge1}\frac{(|t|/\rho)^j}{j+1}<\infty
\]
uniformly. Also $|\phi(z)|\le e^{|z|}/2$, from its elementary
integral representation
$\phi(z)=\int_0^1(1-s)e^{sz}\,ds$.
Apart from a uniform constant, the absolute value of the $r$th
summand of \eqref{hp:Nanalytic} is consequently at most
\[
 \frac{|x_r|^2}{4^r r!(r+1)!}
 \le\frac{M^2|t|^2}{r+1}
                  \left(\frac{|t|}{2\rho}\right)^{2r}.
\]
This is summable. The same reasoning applies on each compact
subset, proving normal convergence and holomorphy. At any fixed
power of $t$ only finitely many $r$ contribute. Expanding the
exponential and $\phi$, then using
\eqref{hp:normal-coefficient}, gives exactly the canonical
coefficients. This proves \eqref{hp:normalTaylor}.
\end{proof}

The factor $\phi$ removes all apparent singularities before any
infinite summation. In particular odd-index kernels cause no
finite-part integral in the analytic formula.

\begin{theorem}[Analytic primitive and translation identity]
\label{hp:analyticH}
Under the same hypothesis, the canonical series
\begin{equation}\label{hp:Hanalytic}
           \mathfrak H_U(x,t)=\sum_{m\ge1}H_m(U(x))\frac{t^m}{m!}
\end{equation}
converges locally uniformly on the smaller domain
\begin{equation}\label{hp:primitivedomain}
                        |x-x_0|+2|t|<R.
\end{equation}
If $F'=U$ locally, then on this domain
\begin{equation}\label{hp:translation}
 \boxed{\displaystyle
 e^{F(x)-F(x-t)}=\mathfrak N_U(x,t)
                               +\partial_x\mathfrak H_U(x,t).}
\end{equation}
The expression on the left is independent of the additive constant
in $F$.
\end{theorem}

\begin{proof}
We provide the convergence justification, since the formal
decomposition alone is insufficient. Let $P=A-N$ denote the
formal generating-series difference. The coefficientwise homotopy
formula can be written as
\begin{equation}\label{hp:Hintegral}
 \mathfrak H_U=\int_0^1\sum_{j\ge1}\sum_{k=0}^{j-1}
 U^{(j-1-k)}(x)(-\partial_x)^k
       \left.(\partial_{u_j}P)\right|_{u_r=\lambda U^{(r)}(x)}
                                      \,d\lambda.
\end{equation}
The factor of $\lambda$ from differentiating $P(\lambda u)$
cancels the $1/\lambda$ in $K^{-1}P=\int_0^1P(\lambda u)d\lambda/\lambda$.
There is therefore no singular endpoint in \eqref{hp:Hintegral}.

Fix a compact subset of \eqref{hp:primitivedomain}. Around each
of its $x$ values choose a disk of radius $\rho>|t|$ whose
closure still satisfies the normal-generator condition with
the same $t$. The raw exponential gives immediately
\[
 |\partial_{u_j}A|\le C\frac{|t|^{j+1}}{(j+1)!}.
\]
Differentiation of \eqref{hp:Nanalytic}, uniformly for
$0\le\lambda\le1$, gives the same bound for
$\partial_{u_j}N$. To see the only potentially different term,
differentiate the $r=j$ summand in its nonsingular form and use
\[
 \left|\frac{d}{dy}\{y^2\phi(cy)\}\right|
   =\left|y\frac{e^{cy}-1}{cy}\right|
   \le |y|e^{|cy|},
\]
with continuous interpretation at $cy=0$. The prefactor
$1/(4^jj!(j+1)!)$ and Cauchy's jet estimate then give
the required bound. For $r>j$ there is just the factor
$t^{j+1}C_j$, and the already summable majorant applies.
These bounds hold on a slightly larger compact set, so Cauchy's
estimate in $x$ yields
\[
 \left|\partial_x^k(\partial_{u_j}P)(\lambda U)\right|
 \le C\frac{k!\,|t|^{j+1}}{(j+1)!\rho^k},\qquad
 |U^{(j-1-k)}(x)|\le M\frac{(j-1-k)!}{\rho^{j-1-k}}.
\]
If necessary shrink $\rho$ slightly to have both bounds with
one common radius still exceeding $|t|$. Since
\[
 \sum_{k=0}^{j-1}\frac{k!(j-1-k)!}{(j+1)!}\le\frac1{j+1},
\]
the $j$th group in \eqref{hp:Hintegral} is bounded by
$CM|t|^{j+1}\rho^{1-j}/(j+1)$, a summable series.
The same bounds prove normal convergence on compact subsets.
Coefficient extraction is justified and reproduces
Theorem~\ref{hp:closed}. Finally the convergent raw series is
$\exp(\int_{x-t}^xU(s)ds)=e^{F(x)-F(x-t)}$.
The coefficientwise equality $A_m=N_m+DH_m$, now between
holomorphic functions, proves \eqref{hp:translation}.
\end{proof}

\begin{corollary}[Gamma and Stieltjes translation formulas]
\label{hp:specialtranslations}
For positive $x$ and $|t|<x/2$,
\begin{align}
 \frac{\Gamma(x)}{\Gamma(x-t)}
     &=\mathfrak N_{\psi}(x,t)
                         +\partial_x\mathfrak H_{\psi}(x,t),
                                                    \label{hp:Gammatrans}\\
 \exp\bigl(F_q(x)-F_q(x-t)\bigr)
     &=\mathfrak N_{-\gamma_q}(x,t)
                         +\partial_x\mathfrak H_{-\gamma_q}(x,t),
                                                    \label{hp:Stieltjestrans}
\end{align}
with $F_q$ as in \eqref{hp:StieltjesU}. These are ordinary
holomorphic identities in $t$ and local holomorphic identities in $x$.
On a positive interval $[a,b]$ with $|t|<a/2$ they may be
integrated termwise and retain the exact endpoint difference
$\mathfrak H_U(b,t)-\mathfrak H_U(a,t)$.
\end{corollary}

\begin{proof}
Both $\psi$ and $\gamma_q$ are holomorphic on the right half-plane.
At a fixed positive $x$, every disk of radius $R<x$ is available.
Take $2|t|<R<x$ in Theorem~\ref{hp:analyticH}, then substitute
the standard primitives. Compactness gives uniformity on $[a,b]$.
\end{proof}

No sharp convergence radius is asserted. The purpose of the domain is
to justify infinite-jet substitution, coefficientwise differentiation,
and anchored integration for actual special functions.

\subsection{Independent finite checks and the remaining question}

The accompanying \texttt{normal\_forms.py} constructs $A_m$ by
its differential recurrence, constructs $N_m$ independently from
the rational beta coefficient formula, and computes $H_m$ in two
ways: the direct canonical reduction and the closed Euler homotopy.
Through weight~20 all results agree, including $A_m=N_m+DH_m$
and the full Euler--boundary identity. There are 106 recorded
exact assertions, including a rejected coefficient corruption.
These checks support implementation correctness; the proofs of
the all-order assertions are the theorems above.

An additional exploratory calculation through weight~24 suggests
that for $2\le m\le24$ every monomial of weight $m-1$ occurs in
$H_m$, with sign
$(-1)^{m-|\boldsymbol a|}$ for its coefficient. We do not use or
claim this as a theorem. A positivity formula for the absolute
coefficients in \eqref{hp:finite-sum} would strengthen the
closed primitive result and explain the observed full support.
